\section{Заключение}

В работе сформулирована вероятностная постановка задачи выбора времени остановки при рекурсивном прогнозировании речи. Речь рассматривается как случайная последовательность на конечном словаре, а рекурсивная модель строит несколько последовательных распределений следующей языковой единицы. По полной последовательности потерь определяется оракульное время остановки, минимизирующее сумму ошибки и стоимости вычислений.

Основная статистическая задача состоит в оценке условного распределения этого оракульного времени по информации, доступной на текущем внутреннем шаге. Правило остановки отделено от модели распределения. Это позволяет учитывать не только вероятность остановки сейчас, но и вероятность того, что оптимальный момент уже был в прошлом или еще находится впереди.

Рассмотрены шесть семейств распределений на \(\{0,\ldots,K\}\): конечное дискретное, сглаженная целевая разметка по потерям, смесь геометрических, дискретизированная смесь экспоненциальных, степенной спад на отрезке и отрицательное биномиальное распределение. Для принятия решения используются пять правил: накопленная вероятность, малая вероятность будущего улучшения, условная интенсивность по предсказанной мере, квантиль и бюджетное правило.

На данных WikiText-103 и ARC-AGI приведены согласованные по протоколу контрольные сводки: перебор стратегий остановки, эмпирическое распределение оракульного \(\tau^\star\) для варианта \texttt{finite\_discrete} и сравнение базового обучения с дополнительным обучением только головы оракула.

\textbf{ВЫВОД.} По результатам экспериментов наиболее устойчиво хорошее качество при контролируемой глубине дают \emph{смесевые} параметризации (\texttt{mixture\_geometric}, \texttt{mixture\_exponential}) в связке с правилом накопленной вероятности или \texttt{hazard}; на языке это согласуется с относительной гладкостью кривой ошибки по внутренним шагам, на ARC --- с необходимостью допускать и ранний выход, и удлинение цепочки на «жёстких'' примерах. Отдельное дообучение головы оракула в основном смещает выбранное правило к более ранней остановке при малом изменении топ-1, что указывает на доминирующую роль базовой репрезентации и умеренную чувствительность постановки к тонкой настройке только классификационного блока распределения момента остановки.

Направления дальнейшей работы: (i) полный языковой перебор по той же схеме оценки; (ii) формальные оценки смещения эмпирического распределения \(\tau^\star\) при переходе от обучающей выборки к контрольной; (iii) постановка вероятностных гарантий устойчивости выбранного правила остановки при смене выборки (в том числе при разных разбиениях на обучение и контроль).
